\begin{frame}{실험 결과: $+$27\% 과대추정 vs $+$5\%}
  \small
  \begin{tabular}{@{}lccc@{}}
    \toprule
     & 최종 $Q(8,\text{오른쪽})$ & 참값(4.720) 대비 & TD 손실 \\
    \midrule
    A. 타겟 없음 & 6.012 & \textbf{$+$27\% 과대추정} & 0.46 \\
    B. 타겟 있음 ($C=200$) & 4.961 & $+$5\% & 1.32 \\
    참값 $Q^{*}$ & 4.720 & --- & --- \\
    \bottomrule
  \end{tabular}
  \normalsize
  \vskip0.5em
  \begin{columns}[T]
    \begin{column}{0.4\textwidth}
      \centering
      \includegraphics[width=\linewidth]{ch09_2_target_net_effect.png}
      \\[0.2em]{\footnotesize (a) $x=8$의 $Q$(오른쪽).
      (b) TD 손실(로그 스케일).}
    \end{column}
    \begin{column}{0.55\textwidth}
      {\footnotesize (마지막 100샘플 평균. 시드 0.
      13상태 환경이라 \textbf{크기}(27\% vs 5\%)는 이 실험의 수치일 뿐,
      \textbf{방향}은 견고하다.)}
    \end{column}
  \end{columns}
\end{frame}
